\originalchapter{电磁波、界面与波导}
\originalsection{横波与能流}
无源均匀真空中对法拉第定律取旋度, 用 $\nabla\times\nabla\times\ee=\nabla(\nabla\cdot\ee)-\nabla^2\ee$, 得 $\nabla^2\ee-c^{-2}\partial_t^2\ee=0$. 磁场同样满足波动方程. 对复表示 $\ee=\Re[\vct E_0e^{i(\vct k\cdot\rr-\omega t)}]$, 源为零给 $\vct k\cdot\vct E_0=0$、$\vct k\cdot\vct B_0=0$, 且 $\vct B_0=(\vct k\times\vct E_0)/\omega$, $\omega=c|\vct k|$. 电、磁、传播方向相互垂直.

\begin{center}
\begin{tikzpicture}[>=Stealth]
\draw[->,thick] (0,0)--(3,0) node[right] {$\vct k$};
\draw[->,thick] (0,0)--(0,1.7) node[above] {$\ee$};
\draw[->,thick] (0,0)--(-1.1,-.8) node[below left] {$\bb$};
\node at (2,1.3) {$\ee\times\bb\parallel\vct k$};
\end{tikzpicture}\par
横波三方向的右手关系, 透视示意, 不表示相同场量的单位或比例.
\end{center}
波阻抗 $Z_0=\sqrt{\mu_0/\eps}$ 给 $|E_0|=Z_0|H_0|$. 时间平均能流 $\langle\vct S\rangle=|E_0|^2\hat{\vct k}/(2Z_0)$. 两个正交电场分量同相为线偏振, 等振幅且相差 $\pi/2$ 为圆偏振, 一般情形为椭圆偏振.

\originalsection{反射与透射}
设两侧为静止、线性、均匀、各向同性的普通无损介质, $\varepsilon_j,\mu_j>0$, 波阻抗 $Z_j=\sqrt{\mu_j/\varepsilon_j}$, 平面界面上没有自由面电荷与面电流. 正入射时电场的连续性为 $E_i+E_r=E_t$, 磁场的连续性为 $(E_i-E_r)/Z_1=E_t/Z_2$. 反射波磁场的负号由传播方向反转产生. 解得
\[
r_E=\frac{Z_2-Z_1}{Z_2+Z_1},\quad t_E=\frac{2Z_2}{Z_2+Z_1},\quad
R=|r_E|^2,\quad T=\frac{Z_1}{Z_2}|t_E|^2,\quad R+T=1.
\]
$|t_E|$ 可以大于1而没有能量增益, 因为能流还含阻抗. 对斜入射, 同一界面上的相位必须一致, 得频率相同、切向波矢相同, 即 $n_1\sin\theta_i=n_2\sin\theta_t$ 与 $\theta_r=\theta_i$. s 偏振的界面波导纳为 $\cos\theta/Z$, p 偏振若以切向电场为振幅, 导纳为 $1/(Z\cos\theta)$. 记相应导纳为 $Y_1,Y_2$, 切向电场连续给 $E_i^\parallel+E_r^\parallel=E_t^\parallel$, 切向磁场连续给 $Y_1(E_i^\parallel-E_r^\parallel)=Y_2E_t^\parallel$. 因而两种偏振都可统一写
\[
r_\parallel=\frac{Y_1-Y_2}{Y_1+Y_2},\quad
 t_\parallel=\frac{2Y_1}{Y_1+Y_2},\quad
 R=|r_\parallel|^2,\quad T=\frac{Y_2}{Y_1}|t_\parallel|^2.
\]
这里假定透射角为实数. p 偏振的总电场透射振幅比为 $t_E=t_\parallel\cos\theta_i/\cos\theta_t$; s 偏振的电场全在界面内, 无此换算. 例如非磁介质 $n_2=\sqrt3n_1$, 取 $\theta_i=\pi/3$, 则 $\theta_t=\pi/6$ 且 $Y_1=Y_2$ 对 p 偏振成立, $r_p=0$, 这就是布儒斯特角 $\tan\theta_B=n_2/n_1$. 振幅是总电场还是切向分量必须先约定, 否则透射系数会差一个余弦因子.

\originalsection{导体与约束模态}
对静止、线性、均匀、各向同性导体取正实近似常数 $\mu,\sigma$ 及正频率 $\omega>0$, $\jj=\sigma\ee$. 对 $e^{-i\omega t}$ 约定, $k^2=\mu\varepsilon\omega^2+i\mu\sigma\omega$. 当 $\sigma\gg\omega\varepsilon$ 时 $k\simeq(1+i)/\delta$, $\delta=\sqrt{2/(\mu\sigma\omega)}$. 波沿 $+z$ 传播时振幅为 $e^{-z/\delta}$, $\delta$ 是趋肤深度, 不是只改变相速度的实参数.

理想导体波导侧壁要求切向电场为零. 沿 $z$ 写 $e^{i(k_zz-\omega t)}$, 横向函数满足 $(\nabla_\perp^2+k_\perp^2)\psi=0$, 而 $k_z^2=(\omega/c)^2-k_\perp^2$. 矩形截面 $a\times b$ 上 TM 模态取 $E_z\propto\sin(m\pi x/a)\sin(n\pi y/b)$, $m,n\ge1$; TE 模态取 $H_z\propto\cos(m\pi x/a)\cos(n\pi y/b)$, $m,n\ge0$ 且不同时为零. 横向本征值 $k_\perp^2=(m\pi/a)^2+(n\pi/b)^2$ 决定截止频率 $\omega_c=ck_\perp$. 这些边界条件不是任取的: TM 模态的 $E_z$ 沿导体侧壁切向, 因而须在边界为零, 是 Dirichlet 条件. 对 TE 模态, 由横向麦克斯韦方程得 $\ee_\perp=-i\omega\mu_0\hat{\vct z}\times\nabla_\perp H_z/k_\perp^2$, 故壁上切向电场零要求 $\partial_nH_z=0$, 是 Neumann 条件 (采用 $e^{ik_zz-i\omega t}$ 约定). 当 $m=n=0$ 时 $H_z$ 为常数; 但对 $k_z\ne0$, $\nabla\cdot\bb=0$ 又要求此常数为零, 所以它没有对应的非零传播场. 单连通空心波导也没有 TEM 模态; 同轴双导体结构可有 TEM 模态.

当 $\omega>\omega_c$ 时 $v_\mathrm{ph}=\omega/k_z>c$, $v_g=\dd\omega/\dd k_z=c^2k_z/\omega<c$, 二者乘积为 $c^2$. 相速度超光速不表示能量或信息超光速. 例如 $a=2b$ 时最低截止模态是 TE$_{10}$, $\omega_c=\pi c/a$; 在 $\omega=2\omega_c$ 下 $k_z=\sqrt3\pi/a$, $v_\mathrm{ph}=2c/\sqrt3$, $v_g=\sqrt3c/2$. TM$_{11}$ 的截止为 $\sqrt5\pi c/a$, 在此频率不能传播. 本章本科部分依据 8.07 第14讲和作业10; 波导为 8.311 补充笔记 \textit{Wave Guides} 支持的进阶扩展.
